% ============================================================================
%  Chapter 0 --- Front Matter (Title Page, Dedications, Abstract)
% ============================================================================

% ── Title Page ────────────────────────────────────────────────────────────────
\begin{titlepage}
\begin{center}

\vspace*{0.5cm}

% --- University & Company logos (replace with actual files) ---
% \includegraphics[height=2cm]{images/esprit-logo.png}
% \hfill
% \includegraphics[height=2cm]{images/huawei-logo.png}

{\large\textsc{ESPRIT --- École Supérieure Privée d'Ingénierie et de Technologies}}\\[4pt]
{\small Department of Computer Science and Telecommunications}

\vspace{1.5cm}

\rule{\linewidth}{0.8pt}\\[0.6cm]
{\LARGE\bfseries Cloud-Native AI Operations Agent\\[6pt]
for CEM--CVM Intelligence}\\[0.3cm]
{\large\itshape Bridging Customer Experience Management and Customer Value Management\\
with Agentic AI on Huawei Cloud Stack}\\[0.4cm]
\rule{\linewidth}{0.8pt}

\vspace{1.5cm}

{\large\textbf{End-of-Studies Project (PFE)}\\[4pt]
Submitted in partial fulfilment of the requirements for the\\
\textbf{Engineering Degree in Computer Science / Telecommunications}}

\vspace{1.5cm}

\begin{tabular}{rl}
\textbf{Student:}     & Souhayl Guenichi \\[4pt]
\textbf{Host Company:}& Huawei Tunisia --- Cloud IT / Sales-Solution \\[4pt]
\textbf{Supervisor:}  & \textit{[Supervisor Name]} \\[4pt]
\textbf{Academic Year:}& 2025 -- 2026 \\
\end{tabular}

\vfill

{\small March 2026}

\end{center}
\end{titlepage}

% ── Dedication (optional) ────────────────────────────────────────────────────
\chapter*{Dedication}
\addcontentsline{toc}{chapter}{Dedication}
\vspace*{2cm}
\begin{center}
\textit{To my family, for their unconditional support throughout this journey.}
\end{center}
\newpage

% ── Acknowledgments ──────────────────────────────────────────────────────────
\chapter*{Acknowledgments}
\addcontentsline{toc}{chapter}{Acknowledgments}

I would like to express my sincere gratitude to:

\begin{itemize}[leftmargin=1.5cm]
  \item My supervisor at Huawei Tunisia, for his guidance on the industrial positioning of this project within the CEM--CVM--ADN ecosystem and for facilitating access to real operator data.
  \item The Huawei Cloud IT / Sales-Solution team in Tunisia, for providing the technical context of SmartCare, Customer Value Management, and the Autonomous Driving Network paradigm.
  \item Tunisie Telecom, for granting access to anonymised production OSS/BSS data that enabled real-world model training and validation.
  \item My academic supervisor at ESPRIT, for the methodological guidance and critical review throughout this project.
  \item The jury members, for the time dedicated to evaluating this work.
\end{itemize}
\newpage

% ── Abstract (English) ───────────────────────────────────────────────────────
\chapter*{Abstract}
\addcontentsline{toc}{chapter}{Abstract}

Modern telecommunications operators generate vast volumes of network performance data (OSS) and business metrics (BSS) across heterogeneous systems. However, the correlation between network degradations and their business impact---customer churn, revenue loss, SLA breaches---remains a manual, fragmented process. Huawei's Customer Experience Management platform (SmartCare) produces rich experience indicators, yet the intelligence layer that translates these signals into actionable business decisions for Customer Value Management (CVM) systems is largely absent or siloed.

This project designs and implements a \textbf{cloud-native AI Operations Agent} that fills this gap. Positioned as the intelligence bridge between CEM (SmartCare) and CVM, the platform ingests real OSS and BSS data from Tunisie Telecom, organises it through a three-layer data lake (Raw $\to$ Processed $\to$ Curated), and applies machine learning models---\texttt{GradientBoostingRegressor} for SLA risk prediction and \texttt{IsolationForest} for network and revenue anomaly detection---to produce actionable insights. An OSS--BSS correlation engine quantifies the statistical relationship between network degradation and revenue impact using Pearson and Spearman coefficients.

The entire platform is containerised with Docker Compose (five microservices: API Gateway, Pipeline Worker, AI Service, PostgreSQL, MinIO) and architecturally mapped to Huawei Cloud Stack (HCS): OBS for object storage, RDS for the relational layer, and ECS/CCE for compute. The system is framed within Huawei's \textbf{Autonomous Driving Network (ADN)} paradigm, demonstrating how an AI agent can autonomously monitor, correlate, and act upon telecom operational and business signals.

Trained on real anonymised data from a live Tunisian operator, the models deliver credible precision, recall, and F1 scores validated against known network events. The project constitutes an industrial-grade reference implementation of O+B (OSS+BSS) convergence at PFE scale.

\vspace{0.5cm}
\noindent\textbf{Keywords:} AI Operations, CEM, CVM, SmartCare, Anomaly Detection, SLA Risk Prediction, OSS--BSS Convergence, Cloud-Native, Docker, Huawei Cloud Stack, Autonomous Driving Network, Tunisie Telecom, Machine Learning, Data Lake.

\newpage

% ── Résumé (French) ──────────────────────────────────────────────────────────
\chapter*{Résumé}
\addcontentsline{toc}{chapter}{Résumé}

Les opérateurs de télécommunications modernes génèrent d'importants volumes de données de performance réseau (OSS) et de métriques métier (BSS) à travers des systèmes hétérogènes. La corrélation entre les dégradations réseau et leur impact métier --- attrition client, perte de revenus, violations de SLA --- reste un processus fragmenté et manuel. La plateforme CEM (Customer Experience Management) de Huawei, SmartCare, produit de riches indicateurs d'expérience, mais la couche d'intelligence qui traduit ces signaux en décisions actionnables pour les systèmes CVM (Customer Value Management) est en grande partie absente.

Ce projet conçoit et implémente un \textbf{Agent AI d'Opérations cloud-natif} qui comble ce vide. Positionné comme le pont d'intelligence entre CEM (SmartCare) et CVM, la plateforme ingère des données OSS et BSS réelles de Tunisie Telecom, les organise via un lac de données à trois couches (Brut $\to$ Traité $\to$ Consolidé), et applique des modèles de machine learning --- \texttt{GradientBoostingRegressor} pour la prédiction du risque SLA et \texttt{IsolationForest} pour la détection d'anomalies réseau et de revenus --- afin de produire des insights actionnables. Un moteur de corrélation OSS--BSS quantifie la relation statistique entre la dégradation réseau et l'impact sur les revenus à l'aide des coefficients de Pearson et Spearman.

L'ensemble de la plateforme est conteneurisée avec Docker Compose (cinq microservices) et architecturalement projetée sur Huawei Cloud Stack (HCS). Le système s'inscrit dans le paradigme du \textbf{Réseau Autonome (ADN)} de Huawei, démontrant comment un agent IA peut surveiller, corréler et agir de manière autonome sur les signaux opérationnels et métier des télécommunications.

\vspace{0.5cm}
\noindent\textbf{Mots-clés :} Agent IA, CEM, CVM, SmartCare, Détection d'anomalies, Prédiction de risque SLA, Convergence OSS--BSS, Cloud-Natif, Docker, Huawei Cloud Stack, Réseau Autonome, Tunisie Telecom, Apprentissage Automatique, Lac de données.

\newpage

% ── Glossary / Acronyms ─────────────────────────────────────────────────────
\chapter*{List of Acronyms}
\addcontentsline{toc}{chapter}{List of Acronyms}

\begin{longtable}{@{}ll@{}}
\toprule
\textbf{Acronym} & \textbf{Definition} \\
\midrule
\endhead
ADN   & Autonomous Driving Network \\
API   & Application Programming Interface \\
APPU  & Average Purchase Per User \\
ARPU  & Average Revenue Per User \\
BSS   & Business Support System \\
CCE   & Cloud Container Engine (Huawei) \\
CEI   & Customer Experience Index \\
CEM   & Customer Experience Management \\
CVM   & Customer Value Management \\
DOU   & Data of Use (monthly data consumption) \\
EAP   & Experience Analytics Platform \\
ECS   & Elastic Cloud Server (Huawei) \\
FTTx  & Fibre to the x (generic fibre access) \\
GBR   & Gradient Boosting Regressor \\
HCS   & Huawei Cloud Stack \\
IF    & Isolation Forest \\
IoT   & Internet of Things \\
KPI   & Key Performance Indicator \\
KQI   & Key Quality Indicator \\
NMS   & Network Management System \\
NOM   & Network Operations Management \\
O+B   & OSS + BSS (convergence) \\
OBS   & Object Storage Service (Huawei) \\
OSS   & Operations Support System \\
RAN   & Radio Access Network \\
RDS   & Relational Database Service (Huawei) \\
REST  & Representational State Transfer \\
RSRP  & Reference Signal Received Power \\
SLA   & Service Level Agreement \\
TT    & Tunisie Telecom \\
VAS   & Value Added Services \\
VPC   & Virtual Private Cloud \\
\bottomrule
\end{longtable}

\newpage
